% GENERATED FILE. Edit canonical lessons, then run npm run generate:books.
% canonical-source-hash: fnv1a64:503d9ba44894a8b2
% canonical-lessons: FR-C305-naufrage, FR-C305-ecraser, FR-C305-rescape, FR-C305-survivre, FR-C305-catastrophe

\chapter{Les rescapés}
\label{ch:fr-les-rescapes}
\hlchaptermodality{305}

\begin{chapteropening}
\textbf{By the end of this chapter:} \emph{I can tell the story of a disaster at sea or in the air: a shipwreck, a crash, survivors, surviving and a catastrophe.}

Complete the last lesson of chapter 305: I can tell the story of a disaster at sea or in the air: a shipwreck, a crash, survivors, surviving and a catastrophe.
\end{chapteropening}

\section[le naufrage]{le naufrage — the shipwreck}

\label{lesson:FR-C305-naufrage}



\begin{quote}
Before the new one: say the French for the buoy, then the French for the lifeguard.
\end{quote}

\subsection*{You'll want to know: le naufrage}
\textbf{le naufrage} — \textquotedblleft{}the shipwreck\textquotedblright{}.

\begin{quote}
Le naufrage a eu lieu près de la côte.
\end{quote}

\begin{quote}
\emph{The shipwreck happened near the coast.}
\end{quote}

\emph{Faire naufrage} — to be shipwrecked.

\begin{cousinweb}
\emph{naufrage} is a learned borrowing from Latin \emph{naufragium} \textquotedblleft{}a shipwreck\textquotedblright{} (EtymDB), from \emph{nāvis} \textquotedblleft{}ship\textquotedblright{} and \emph{frangere} \textquotedblleft{}to break\textquotedblright{}.
\end{cousinweb}

\begin{grammarlens}[title={What you've built}]
One more word for telling, explaining and describing at B1.
\end{grammarlens}

\subsection*{Your turn}
\begin{itemize}
\raggedright
  \item \emph{Say it:} \emph{le naufrage}
  \item \emph{Say it:} \emph{le naufrage}, once more
  \item \emph{Say it:} a sentence of your own with \emph{le naufrage}
  \item \emph{Recall:} say the French for the buoy, then the French for the lifeguard, then say \emph{le naufrage} again
\end{itemize}

\begin{tcolorbox}[breakable,colback=teal!4,colframe=teal!35!black,title={Before you move on}]
What does le naufrage mean? (\textquotedblleft{}The shipwreck\textquotedblright{}.) Say it once more.
\end{tcolorbox}
\section[écraser]{écraser — to crash (as s'écraser); to crush}

\label{lesson:FR-C305-ecraser}



\begin{quote}
Before the new one: say the French for the lifeguard, then the French for the shipwreck.
\end{quote}

\subsection*{You'll want to know: écraser}
\textbf{écraser} — \textquotedblleft{}to crash (as s'écraser); to crush\textquotedblright{}.

\begin{quote}
Un petit avion s'est écrasé dans la forêt.
\end{quote}

\begin{quote}
\emph{A small plane crashed in the forest.}
\end{quote}

Reflexive for a plane that crashes: \emph{s'écraser}. Without \emph{se} it is common too, to crush or to mash: \emph{écraser des pommes de terre}.

\begin{cousinweb}
\emph{écraser} comes from Middle French \emph{escraser} \textquotedblleft{}to squash\textquotedblright{} (EtymDB), usually traced to a Scandinavian word, Old Norse \emph{krasa} \textquotedblleft{}to shatter\textquotedblright{}.
\end{cousinweb}

\begin{grammarlens}[title={What you've built}]
One more word for telling, explaining and describing at B1.
\end{grammarlens}

\subsection*{Your turn}
\begin{itemize}
\raggedright
  \item \emph{Say it:} \emph{écraser}
  \item \emph{Say it:} \emph{écraser}, once more
  \item \emph{Say it:} a sentence of your own with \emph{écraser}
  \item \emph{Recall:} say the French for the lifeguard, then the French for the shipwreck, then say \emph{écraser} again
\end{itemize}

\begin{tcolorbox}[breakable,colback=teal!4,colframe=teal!35!black,title={Before you move on}]
What does écraser mean? (\textquotedblleft{}To crash (as s'écraser); to crush\textquotedblright{}.) Say it once more.
\end{tcolorbox}
\section[le rescapé]{le rescapé — the survivor (of a disaster)}

\label{lesson:FR-C305-rescape}



\begin{quote}
Before the new one: say the French for the shipwreck, then the French for to crash.
\end{quote}

\subsection*{You'll want to know: le rescapé}
\textbf{le rescapé} — \textquotedblleft{}the survivor (of a disaster)\textquotedblright{}.

\begin{quote}
Les rescapés ont raconté leur nuit en mer.
\end{quote}

\begin{quote}
\emph{The survivors told the story of their night at sea.}
\end{quote}

Someone saved from a disaster: \emph{les rescapés du naufrage}.

\begin{cousinweb}
\emph{rescapé} is usually explained as a northern dialect form of \emph{réchappé} \textquotedblleft{}escaped\textquotedblright{}, built on \emph{échapper} \textquotedblleft{}to escape\textquotedblright{}, which comes from Old French \emph{eschaper} (EtymDB).
\end{cousinweb}

\begin{grammarlens}[title={What you've built}]
One more word for telling, explaining and describing at B1.
\end{grammarlens}

\subsection*{Your turn}
\begin{itemize}
\raggedright
  \item \emph{Say it:} \emph{le rescapé}
  \item \emph{Say it:} \emph{le rescapé}, once more
  \item \emph{Say it:} a sentence of your own with \emph{le rescapé}
  \item \emph{Recall:} say the French for the shipwreck, then the French for to crash, then say \emph{le rescapé} again
\end{itemize}

\begin{tcolorbox}[breakable,colback=teal!4,colframe=teal!35!black,title={Before you move on}]
What does le rescapé mean? (\textquotedblleft{}The survivor (of a disaster)\textquotedblright{}.) Say it once more.
\end{tcolorbox}
\section[survivre]{survivre — to survive}

\label{lesson:FR-C305-survivre}



\begin{quote}
Before the new one: say the French for to crash, then the French for the survivor.
\end{quote}

\subsection*{You'll want to know: survivre}
\textbf{survivre} — \textquotedblleft{}to survive\textquotedblright{}.

\begin{quote}
Tous les passagers ont survécu.
\end{quote}

\begin{quote}
\emph{All the passengers survived.}
\end{quote}

\emph{Survivre à} — to survive something: \emph{survivre à un accident}.

\begin{cousinweb}
\emph{survivre} is inherited from Late Latin \emph{supervīvere} \textquotedblleft{}to outlive\textquotedblright{} (EtymDB), from \emph{super} \textquotedblleft{}over\textquotedblright{} and \emph{vīvere} \textquotedblleft{}to live\textquotedblright{}.
\end{cousinweb}

\begin{grammarlens}[title={What you've built}]
One more word for telling, explaining and describing at B1.
\end{grammarlens}

\subsection*{Your turn}
\begin{itemize}
\raggedright
  \item \emph{Say it:} \emph{survivre}
  \item \emph{Say it:} \emph{survivre}, once more
  \item \emph{Say it:} a sentence of your own with \emph{survivre}
  \item \emph{Recall:} say the French for to crash, then the French for the survivor, then say \emph{survivre} again
\end{itemize}

\begin{tcolorbox}[breakable,colback=teal!4,colframe=teal!35!black,title={Before you move on}]
What does survivre mean? (\textquotedblleft{}To survive\textquotedblright{}.) Say it once more.
\end{tcolorbox}
\section[la catastrophe]{la catastrophe — the catastrophe, the disaster}

\label{lesson:FR-C305-catastrophe}



\begin{quote}
Before the new one: say the French for the survivor, then the French for to survive.
\end{quote}

\subsection*{You'll want to know: la catastrophe}
\textbf{la catastrophe} — \textquotedblleft{}the catastrophe, the disaster\textquotedblright{}.

\begin{quote}
Le tremblement de terre a été une catastrophe.
\end{quote}

\begin{quote}
\emph{The earthquake was a catastrophe.}
\end{quote}

Informal: \emph{c'est la catastrophe !} — it is a disaster!

\begin{cousinweb}
\emph{catastrophe} was borrowed from Latin \emph{catastropha}, from Greek \emph{katastroph\'{\={e}}} \textquotedblleft{}an overturning, a sudden end\textquotedblright{} (EtymDB).
\end{cousinweb}

\begin{grammarlens}[title={What you've built}]
One more word for telling, explaining and describing at B1.
\end{grammarlens}

\subsection*{Your turn}
\begin{itemize}
\raggedright
  \item \emph{Say it:} \emph{la catastrophe}
  \item \emph{Say it:} \emph{la catastrophe}, once more
  \item \emph{Say it:} a sentence of your own with \emph{la catastrophe}
  \item \emph{Recall:} say the French for the survivor, then the French for to survive, then say \emph{la catastrophe} again
\end{itemize}

\begin{tcolorbox}[breakable,colback=teal!4,colframe=teal!35!black,title={Before you move on}]
What does la catastrophe mean? (\textquotedblleft{}The catastrophe, the disaster\textquotedblright{}.) Say it once more.
\end{tcolorbox}
